\section{Discussion}
\label{sec:discussion}

\paragraph{Two channels, one policy.}
The concern channel is a better error detector than the answer channel: a harness that reads only the answer discards most of what the agent noticed, in every vendor's models. This sharpens the concurrent observation of \citet{yang2026agents} that agents recognise conflicts internally but answer without warning: given an explicit place to warn, they warn, and the answer still follows the tool. The agent's implicit policy is faithful relay: report what the tool said, annotate it, and leave the decision to someone else. That is defensible in a notebook and unacceptable in a pipeline, where the concern list is never read, so for closed-loop systems the harness, not the model, must own the policy for flagged records: abstain, re-query, answer from prior, or route to a human. A one-line policy already shifts seven of ten models, and what the models do on their own is repair, not refuse.

\paragraph{What plausibility checking is made of.}
The sweep, the cue labels and the record-format switch agree that detection has two sources---a contradiction inside the record and recall of the specific entity---and that a general model of what a protein or molecule can be plays almost no role. Correction has a narrower source still: a route from something the agent trusts to a replacement value. A partner field at a known ratio, a sequence to recount, a formula to sum, or a sign to flip all qualify; a remembered ``typically $\sim$364 residues'' does not, and neither does the knowledge that a count cannot be negative when nothing says what it should be. The chemistry results add a threshold to this picture. Models that recall a formula detect a 20\% discrepancy and still defer to the tool, while overriding a 100-fold one; the override is gated not on whether the model has evidence but on whether the evidence is decisive. This is a reasonable policy for a system that cannot tell a stale memory from a bad record, and it explains why capability alone does not close the gap: a stronger model recalls more and so flags more, but it does not override more unless the record gives it something exact to override with.

\paragraph{Design consequences.}
A tool that returns redundant fields gives the agent something to check, and it checks: the full UniProt entry takes impossible propagation to 0--11\% and the PubChem formula to zero on every numeric field, for the same models that relayed 33--62\% from the summary. A tool that returns a bare scalar leaves the agent to its memory, which it will cite and then ignore. Where a record cannot carry redundancy---an identifier, a name---no amount of model capability helps, and the harness has to supply the check (a reverse lookup, a second source) or treat the field as untrusted. The verification instruction is cheap and effective for impossible values on most models, but it works by raising the flag rate, which the harness must then act on; it does little for plausible values, which is where silent errors in real pipelines will live.

\paragraph{Limitations.}
The tools are mocks; real clients return richer records and real errors, and the full-record experiment covers three models in one domain and two in the other. The flag heuristic is keyword-based and the cue classifier rule-based; both are deterministic and released, with a blind validation set, but the chemistry vocabulary was written after the chemistry runs and has not been scored against human labels. The direction split applies to numeric fields only. Default reasoning settings were used throughout; effort levels above the default are untested, and the pilot ablation without extended thinking covers one model. Cue labels record stated reasons, not necessarily what caused a flag. Ten models is enough to see that the pattern is not a vendor artefact, not enough to say which training choices produce the OpenAI models' low recall or the Anthropic models' resistance to the verification sentence.
